\documentclass[10pt,a4paper]{amsart}
\usepackage[margin=19mm]{geometry}
\usepackage{amsmath,amssymb,mathtools}
\usepackage[scr=rsfs,cal=euler]{mathalfa}
\usepackage{xfrac}
\usepackage[unicode,hidelinks,bookmarks=false]{hyperref}
\newcommand{\ie}{i.e.\ }
\newcommand{\cf}{cf.\ }
\newcommand{\resp}{respectively\ }
\newcommand{\Subsection}{\subsection}
\providecommand{\nobreakdash}{\nobreak\hbox{-}}
\setlength{\emergencystretch}{2em}
\multlinegap0pt
\usepackage{bm}
\usepackage{bbm}
\usepackage{dsfont}
\usepackage{xspace}
\usepackage{cancel}
\usepackage{mathtools}
\usepackage{amsthm}
\makeatletter
\@ifundefined{theorem}{\newtheorem{theorem}{Theorem}}{}
\@ifundefined{lemma}{\newtheorem{lemma}[theorem]{Lemma}}{}
\makeatother
\title[Fenchel-Willmore inequality for submanifolds in manifolds wi]{Fenchel-Willmore inequality for submanifolds in manifolds with non-negative $k$-Ricci curvature}
\author{Meng Ji, Kwok-Kun Kwong}
\date{Source: arXiv:2507.07655v2 (CC BY 4.0)}
\begin{document}
\maketitle

\noindent\emph{Source and licence.} Fenchel-Willmore inequality for submanifolds in manifolds with non-negative $k$-Ricci curvature. Version arXiv:2507.07655v2 (CC BY 4.0), distributed under the Creative Commons Attribution 4.0 International licence, as recorded in the arXiv licence metadata for this exact version and shown on the abstract page https://arxiv.org/abs/2507.07655v2. Full rights notice: licenses/arxiv-2507.07655-CC-BY-4.0.txt.\par\smallskip

\noindent\textit{Editorial scope.} Counted unit: the Lemma whose source label is \texttt{lem xyt=1} in the article (source0.tex lines 568--573, source label \texttt{lem xyt=1}), delivered together with the article's own complete original proof of it (source0.tex lines 575--710, one \texttt{proof} environment of 136 lines). No mathematics is written, repaired or re-derived: statement and proof are the article's own text line for line. Required context retained uncounted: thm general fenchel willmore (lines 81--87); thm general fenchel willmore' (lines 148--158); lem 2.2'' (lines 185--191); lem 2.4'' (lines 240--248); ineq est I0 I+ (lines 362--370); ineq I+ (lines 389--394); ineq I+2 (lines 409--415); ineq I+3 (lines 424--430); ineq I+4 (lines 434--437); lem no focal (lines 528--531). Every definition, display and label of those lines is retained unchanged. Omitted from this package and not redistributed: 1--80, 88--147, 159--184, 192--239, 249--361, 371--388, 395--408, 416--423, 431--433, 438--527, 532--567, 574--574, 711--965. Nothing inside a retained excerpt is omitted or altered: every retained body is delivered exactly as the article's lines are written, so no omission receipt of any kind accompanies this package. Citations: the retained excerpts carry no citation command at all, so no bibliography entry of the article is transcribed and no reference is redistributed. Reference adaptation: none was needed and none was made; every reference inside the delivered unit resolves inside it, so no unresolved reference or citation is produced. Printed slips kept literal: no printed word, symbol, hypothesis or number of the counted statement or of its proof was altered. Layout and class: the article's own class is \texttt{amsart} with packages including geometry, amsmath, amssymb, amsthm, amsfonts, amscd, bm, cancel, amsrefs, xcolor, hyperref; the wrapper is the lane's pinned \texttt{amsart} editorial wrapper at 10pt on A4, which restates the theorem-like environments the retained text needs and adds the article's own macro definitions that the retained text needs (none), with extra wrapper packages (bm, bbm, dsfont, xspace, cancel, mathtools). The delivered numbering is the unit's own. Assets: no retained span contains a figure environment, a graphics-inclusion command, a PDF-page inclusion, a table or a TikZ picture; no image file is copied into this package, the package declares no asset dependency and no image file is redistributed with it. Licence: version arXiv:2507.07655v2 of this article is distributed under the Creative Commons Attribution 4.0 International licence, as recorded in the arXiv licence metadata for this exact version and shown on the abstract page https://arxiv.org/abs/2507.07655v2. A full rights notice is delivered at licenses/arxiv-2507.07655-CC-BY-4.0.txt. Characters: every retained excerpt is pure ASCII, so the delivered engine, XeLaTeX with its default Latin Modern text font, reports no missing character.
\begin{theorem}\label{thm general fenchel willmore}
Let $n, m \in \mathbb{N}$, and $(M, g)$ be a complete non-compact Riemannian manifold of dimension $n+m$ with nonnegative sectional curvature and asymptotic volume ratio $\theta>0$. Suppose $\Sigma$ is a closed $n$-dimensional submanifold immersed in $M$. Then
\begin{equation}\label{ineq fenchel willmore}
\int_{\Sigma}|\sigma|^n \ge \theta|\mathbb S^n|.
\end{equation}
If $m=1$, the curvature assumption may be relaxed to require only non-negative Ricci curvature.
\end{theorem}

\begin{theorem}\label{thm general fenchel willmore'}
The conclusion of Theorem \ref{thm general fenchel willmore} still holds under the weaker assumption that the $k$-Ricci curvature is non-negative, where
\begin{equation}\label{eq k}
k=
\begin{cases}
\min(n, m-1), & \text{if } m > 1, \\
n, & \text{if } m = 1.
\end{cases}
\end{equation}
Note that in both cases, the Ricci curvature of $M$ is non-negative and $\theta$ is still well-defined.
\end{theorem}

\begin{lemma}\label{lem 2.2''}
For every $0 <\alpha<1$ and $r>0$,
$$
\{p \in M: \alpha r<d(x, p)<r \text { for all } x \in \Sigma\}
$$
is contained in $\left\{\Phi_r(x, z):(x, z) \in A_r \text { and }|z|>\alpha\right\}$.
\end{lemma}

\begin{lemma}\label{lem 2.4''}
Assume that $(x, z) \in A_r$. Then the function
$s \mapsto \frac{|\det D\Phi_s(x, z)|}{s^m(1-s\langle \sigma(x), z\rangle)^n}$
is monotone decreasing and
\begin{align*}
|\det D\Phi_s(x, z)|\le {s^m(1-s\langle \sigma, z\rangle)^n}.
\end{align*}
for $s \in(0, r)$.
\end{lemma}

\begin{equation}\label{ineq est I0 I+}
\begin{aligned}
& |\{p \in M: \alpha r<d(x, p)<r \text { for all } x \in \Sigma\} | \\
\le & \int_{\Sigma} \int_{Z_\alpha}\left|\det D\Phi_r(x, z)\right| 1_{A_r}(x, z) dz d \operatorname{vol}_{\Sigma} \\
=& \int_{\Sigma^0} \int_{Z_\alpha}\left|\det D\Phi_r(x, z)\right| 1_{A_r}(x, z) dz d \operatorname{vol}_{\Sigma}
+ \int_{\Sigma^+} \int_{Z_\alpha}\left|\det D\Phi_r(x, z)\right| 1_{A_r}(x, z) dz d \operatorname{vol}_{\Sigma}\\
=& I^0(\alpha, r)+I^+(\alpha, r),
\end{aligned}
\end{equation}

\begin{equation}\label{ineq I+}
\begin{aligned}
I^+(\alpha, r) & \le \int_{\Sigma^+} \int_{-1}^{\frac{1}{r|\sigma(x)|}} \int_{Y_{\alpha, x, t}}\left|\det D \Phi_r(x, y, t)\right| 1_{A_r}(x, y, t) d y d t d \operatorname{vol}_{\Sigma} (x)\\
& \le \int_{\Sigma^+} \int_{-1}^{\frac{1}{r|\sigma|}} \int_{Y_{\alpha, x, t}} r^m(1-r|\sigma(x)| t)^n d y d t d \operatorname{vol}_{\Sigma},
\end{aligned}
\end{equation}

\begin{equation}\label{ineq I+2}
\begin{aligned}
& \int_{-1}^{\frac{1}{r|\sigma|}} \int_{Y_{\alpha, x, t}} r^m(1-r|\sigma| t)^n d y d t\\
=& (m-1)\left|\mathbb{B}^{m-1}\right|r^m
\int_{-1}^{\frac{1}{r|\sigma|}} \left(\left(1-t^2\right)^{\frac{m-3}{2}}(1-\alpha)+O\left((1-\alpha)^2\right)\right)(1-r|\sigma| t)^n d t.
\end{aligned}
\end{equation}

\begin{equation}\label{ineq I+3}
\begin{aligned}
\int_{-1}^{\frac{1}{r|\sigma|}} \left(1-t^2\right)^{\frac{m-3}{2}} (1-r|\sigma| t)^n d t
=& r^n|\sigma|^n \int_{-1}^0(-t)^n\left(1-t^2\right)^{\frac{m-3}{2}} d t+O\left(r^{n-1}\right)\\
=& \frac{\Gamma\left(\frac{n+1}{2}\right) \Gamma\left(\frac{m-1}{2}\right)}{2\Gamma\left(\frac{n+m}{2}\right)} r^n|\sigma|^n+O\left(r^{n-1}\right).
\end{aligned}
\end{equation}

\begin{equation}\label{ineq I+4}
\int_{-1}^{\frac{1}{r|\sigma|}}(1-r|\sigma| t)^n d t=
\frac{(1+r|\sigma|)^{n+1}}{(n+1) r|\sigma|}=O(r^n).
\end{equation}

\begin{lemma}\label{lem no focal}
For all $x \in \Sigma^+$, $z \in {T}_x^{\perp} \Sigma$ such that $|z|=1$ and $\langle z, \sigma(x)\rangle \le 0$,
the geodesic $\gamma(s)=\exp_x\left(sz\right)$ minimizes the distance from $\Sigma$ for all $s \ge0$. In particular, no focal point of $\Sigma$ occurs along $\gamma(s)$ for $s>0$ and the normal exponential map when restricted to $\{(x, z)\in T^\perp \Sigma^+: \langle z, \sigma(x)\rangle \le 0\}$ is injective, and hence is a diffeomorphism onto its image.
\end{lemma}

\begin{lemma}\label{lem xyt=1}
For every $r>0, x \in \Sigma^+, z \in {T}_x^{\perp} \Sigma$ satisfying $|z|=1$ and $\langle z, \sigma(x)\rangle \le0$, we have
$$
\left|\det D \Phi_r(x, z)\right| \ge r^m\left(1-r\langle z, \sigma(x) \rangle \right)^n.
$$
\end{lemma}

\begin{proof}

Assume on the contrary that there exists $x_{0}\in \Sigma^+$, $z_{0}\in {T}_{x}^{\perp}\Sigma$ that satisfy $|z_{0}|=1$ and $\langle z_0, \sigma(x_0)\rangle \le0$, such that
$$
\left|\det D \Phi_{r_{0}}(x_{0}, z_0)\right| < r_{0}^m\left(1-r_0\langle z_0, \sigma(x_0)\rangle \right)^n
$$
for some $r_{0}>0$. Since this is an open condition, we can without loss of generality assume that $y_0\ne0$.

By continuity, there exist $\varepsilon\in(0, 1)$ and a neighbourhood $V$ of $(x_{0}, z_0)$ in $T^{\perp}\Sigma$ such that
$$
\left|\det D\Phi_{r_{0}}(x, z)\right|
< (1-\varepsilon) r_{0}^{m} \bigl(1-r_{0}\langle z, \sigma(x)\rangle \bigr)^{n}
\quad\text{on } V.
$$
Lemma \ref{lem no focal} and Lemma \ref{lem 2.4''} then imply that for every $r>r_{0}$,
$$
\left|\det D\Phi_{r}(x, z)\right|
< (1-\varepsilon) r^{m} \bigl(1-r \langle z, \sigma(x)\rangle \bigr)^{n}
\quad\text{on } V\cap A_{r}.
$$

Now, assume that $m\ge 3$ and write $(x, z)=(x, y, t)$ as before.
For $\alpha\in(0, 1)$ and each $(x, t)$, define
$$
Y_{\alpha, x, t}
=\bigl\{ y\in\widetilde T_{x}^{\perp}\Sigma:
\alpha^{2}<|y|^{2}+t^{2}<1\bigr\},
$$
and write $Y_{\alpha}=Y_{\alpha, x, t}$.

Consequently, by applying Lemma \ref{lem 2.2''}, and following the reasoning in \eqref{ineq est I0 I+}, \eqref{ineq I+}, \eqref{ineq I+2}, \eqref{ineq I+3} and \eqref{ineq I+4}, we have
\begin{equation}\label{ineq est equality case}
\begin{aligned}
& |\{p \in M: \alpha r<d(x, p)<r \text { for all } x \in \Sigma\}| \\
\le& \int_{\Sigma^+} \int_{-1}^{\frac{1}{r|\sigma|}} \int_{Y_\alpha}
\left|\det D \Phi_r(x, y, t)\right| 1_{A_r}(x, y, t) d y d t d \mathrm{vol}_{\Sigma} \\
\le& \int_{\Sigma^+} \int_{-1}^{\frac{1}{r|\sigma|}} \int_{Y_\alpha}
\big(1-\varepsilon 1_{V}(x, y, t)\big) r^m\left(1-r|\sigma| t\right)^n d y d t d \mathrm{vol}_{\Sigma}+I^{0}(\alpha, r)\\
\le&
(1-\alpha)\frac{\Gamma\left(\frac{n+1}{2}\right) \Gamma\left(\frac{m-1}{2}\right)}{2 \Gamma\left(\frac{n+m}{2}\right)} r^{m+n}\int_{\Sigma^+}|\sigma|^n\\
& - \varepsilon \int_{\Sigma^+} \int_{-1}^{\frac{1}{r|\sigma|}} \int_{Y_\alpha}
1_{V}(x, y, t) r^m\left(1-r|\sigma| t\right)^n d y d t d \mathrm{vol}_{\Sigma}\\
& +O(1-\alpha)O\left(r^{m+n-1}\right) +O((1-\alpha)^2)O(r^{m+n})\\
=&:J(\alpha, r)-\varepsilon I(\alpha, r)+\mathrm{remainder}
\end{aligned}
\end{equation}
for all $r>r_{0}$ and $\alpha$ close to $1$.

We now claim that $\displaystyle \liminf_{\alpha \to 1^{-}} \liminf_{r \to \infty} \frac{1}{1-\alpha}\cdot\frac{1}{r^{n+m}} I(\alpha, r)$ is positive. This is straightforward but slightly tedious.

As $V$ is open, for $\alpha$ sufficiently close to 1 the set $V \cap Y_{\alpha, x_0, t_0}$ contains, in polar coordinates on $\widetilde{T}_{x_0}^{\perp} \Sigma$, an open set
$\{(\rho, \theta):\ \alpha^{2}-t_{0}^{2}<\rho^{2}<1-t_{0}^{2}, \ \theta\in\mathcal O\}$,
where $\mathcal O\subset\mathbb S^{m-2}$ is an open neighbourhood of $\frac{y_{0}}{|y_{0}|}$ independent of $\alpha$.

Hence, for $\alpha$ close to $1$,
$$
|V\cap Y_{\alpha, x_{0}, t_{0}}|
\;\ge\;
\frac{|\mathcal O|}{m-1}\Bigl[(1-t_{0}^{2})^{\frac{m-1}{2}}-(\alpha^{2}-t_{0}^{2})^{\frac{m-1}{2}}\Bigr]
\ge2 \delta_1(1-\alpha)
$$
for some $\delta_{1}>0$ independent of $\alpha$. By shrinking $V$ if necessary, we may assume $|V\cap Y_{\alpha, x, t}|\ge\delta_{1}(1-\alpha)$ for all $(x, y, t)\in V$ and $\alpha$ close to $1$.

Consequently, for $\alpha$ near $1$,
\begin{equation}\label{ineq lower bound 2}
\frac{1}{1-\alpha}\cdot \frac{1}{r^{n+m}} I(\alpha, r) \ge \frac{\delta_{1}}{r^{n}}\int_{B_{\rho}(x_{0})} \int_{T_{x, r}\cap V} \bigl(1-r|\sigma| t\bigr)^{n}dt\, d\operatorname{vol}_{\Sigma},
\end{equation}
for some $\rho>0$, where $T_{x, r} =\{t\in \mathbb R: -1<t<\tfrac{1}{r|\sigma(x)|}\}$.

Because the integrand of the integral on the RHS in \eqref{ineq lower bound 2} is
decreasing in $t$, we have
\begin{equation}\label{ineq TV}
\begin{aligned}
\int_{T_{x_{0}, r} \cap V}
\bigl(1-r|\sigma(x_{0})| t\bigr)^{n} dt
\ge& \int_{ t_0 -\delta_{2}}^{ t_0 } \bigl(1-r|\sigma(x_{0})| t\bigr)^{n} dt \\
\ge& \int_{ -\delta_{2}}^{0} \bigl(1-r|\sigma(x_{0})| t\bigr)^{n} dt \\
=& \frac{1}{r|\sigma(x_{0})|(n+1)}
\left[\bigl(1+r|\sigma(x_{0})|\delta_{2}\bigr)^{n+1}-1\right]\\
\ge& \frac{\delta_2^{n+1}r^n|\sigma(x_0)|^{n}}{n+1},
\end{aligned}
\end{equation}
where $0<\delta_{2}<1$ is chosen so that $V \cap \{x=x_{0}, y=y_{0}\}$ contains the interval
$(t_{0}-\delta_{2}, t_{0})$.
By continuity,
$$
\int_{T_{x, r} \cap V}
\bigl(1-r|\sigma| t\bigr)^{n} dt
\ge\frac{\delta_{2}^{\, n+1} r^{n}}{2(n+1)} |\sigma(x_{0})|^{n}
$$
for $x$ near $x_{0}$.

Combining this with \eqref{ineq lower bound 2}, we find that
$$
\liminf_{\alpha \rightarrow 1^{-}} \liminf_{r \rightarrow \infty}
\frac{1}{1-\alpha}\cdot \frac{1}{r^{n+m}} I(\alpha, r) > 0.
$$

Putting this into \eqref{ineq est equality case} and proceeding as in the proof of Theorem \ref{thm general fenchel willmore'}, we conclude that
$$
\theta|\mathbb S^{n}|< \int_{\Sigma} |\sigma|^{n},
$$
which is a contradiction.

If $m=2$, note that the condition in Theorem \ref{thm general fenchel willmore'} is that the sectional curvature of $M$ is non-negative. By taking the product of $M$ with $\mathbb{R}$, we can view $\Sigma$ as a codimension $3$ submanifold. Observe that the product manifold $M \times \mathbb{R}$ still has nonnegative sectional curvature, while both $\theta$ and $\int_{\Sigma}|\sigma|^n$ remain invariant. The preceding argument can still be applied, which yields a contradiction.

When $m = 1$, the previous argument requires a slight modification. Observe that $z_0 = -\frac{\sigma(x_0)}{|\sigma(x_0)|}$, and hence, under our identification, $t_0 = -1$. We choose a sufficiently small neighborhood $V$ of $(x_0, t_0)$ as before, with the additional property that for every $(x, t)\in V$, and for all $\alpha$ sufficiently close to $1$ and $r$ sufficiently large, one has $T_{\alpha, r, x}\cap V_x \supset (-1, -\alpha)$. Here $T_{\alpha, r, x}:= \{ t \in \mathbb{R}: -1 < t < \tfrac{1}{r|\sigma(x)|}, \ \alpha < |t| < 1 \}$ and $V_x:= \{ t \in \mathbb{R}: (x, t)\in V \}$.

We replace \eqref{ineq est equality case} with
\begin{equation}\label{ineq est equality case'}
\begin{aligned}
& | \{p \in M: \alpha r<d(x, p)<r \text { for all } x \in \Sigma\} | \\
\le & \int_{\Sigma^{+}} \int_{T_{\alpha, r, x}} \left|\det D \Phi_r(x, t)\right| 1_{A_r}(x, t) d t d \operatorname{vol}_{\Sigma} \\
\le & \int_{\Sigma^{+}} \int_{T_{\alpha, r, x}} \left(1-\varepsilon 1_V(x, t)\right) r(1-r|\sigma| t)^n d t d \operatorname{vol}_{\Sigma}+I^0(\alpha, r)\\
=& \int_{\Sigma^{+}} \int_{T_{\alpha, r, x}} r(1-r|\sigma| t)^n d t d \operatorname{vol}_{\Sigma}
-\varepsilon\int_{\Sigma^{+}} \int_{T_{\alpha, r, x}\cap V_x} r(1-r|\sigma| t)^n d t d \operatorname{vol}_{\Sigma}
+I^0(\alpha, r).
\end{aligned}
\end{equation}
We replace \eqref{ineq TV} with
\begin{equation*}
\begin{aligned}
\int_{T_{\alpha, r, x_0 }\cap V_{x_0}}r \left(1-r\left|\sigma\left(x_0\right)\right| t\right)^n d t
\ge& r \int_{-1}^{-\alpha}\left(1-r\left|\sigma\left(x_0\right)\right| t\right)^n d t\\
=& \frac{\left(1+r\left|\sigma\left(x_0\right)\right|\right)^{n+1}-\left(1+\alpha r\left|\sigma\left(x_0\right)\right|\right)^{n+1}}{(n+1)\left|\sigma\left(x_0\right)\right|}\\
=& {r\left(1+r\left|\sigma\left(x_0\right)\right|\right)^n}(1-\alpha)+O(r^{n+1}(1-\alpha)^2)
\end{aligned}
\end{equation*}
if $\alpha$ is close enough to $1$ and $r$ sufficiently large.

By continuity, it is then easy to see that
$$
\liminf_{\alpha \rightarrow 1^{-}} \liminf_{r \rightarrow \infty} \frac{1}{1-\alpha} \cdot \frac{1}{r^{n+1}}\int_{\Sigma^{+}} \int_{T_{\alpha, r, x}\cap V_x} r(1-r|\sigma| t)^n d t d \operatorname{vol}_{\Sigma}>0.
$$
In view of \eqref{ineq est equality case'}, the same reasoning as in the previous argument applies, which leads to a contradiction.
\end{proof}

\end{document}
